problem:  
Let \((M^{2n+1}, D, J)\) be a compact strictly pseudoconvex CR manifold, where \(D = \ker(\eta)\) for some contact form \(\eta\), and \(J\) is the CR almost complex structure on \(D\). Denote by \(\mathcal{S}(D,J)\) the set of all contact forms \(\tilde{\eta} = e^{f}\eta\) whose Reeb vector fields \(\tilde{\xi}\) preserve \(J\) (that is, \((D,J,\tilde{\eta})\) remain of K-contact type). Each such \(\tilde{\eta}\) determines a corresponding compatible metric  
\[
g_{\tilde{\eta}}(X,Y) = d\tilde{\eta}(X, JY) + \tilde{\eta}(X)\tilde{\eta}(Y),
\]
whose scalar curvature we denote \(s_{\tilde{\eta}}\).  

Consider the total scalar curvature functional restricted to this CR deformation class:
\[
\mathcal{S}(\tilde{\eta}) = \int_M s_{\tilde{\eta}}\, d\mu_{\tilde{\eta}}.
\]
It is known that when the CR structure \((D, J)\) comes from a Sasakian structure, the **transverse scalar curvature** variation leads to extremal Sasakian metrics (via the Sasaki–Futaki invariant).  
However, the global variational geometry of the total *Riemannian* scalar curvature functional over \(\mathcal{S}(D,J)\) remains unclear.

**Question:**  
For a fixed strictly pseudoconvex CR structure \((D, J)\), determine whether critical points of the total scalar curvature functional \(\mathcal{S}(\tilde{\eta})\), under volume normalization and variations of \(\tilde{\eta}\in \mathcal{S}(D,J)\), must necessarily correspond to Sasakian metrics.  

In other terms, does the criticality of the *Riemannian* scalar curvature functional over all K-contact deformations of a fixed CR structure select Sasakian structures as the unique extremals?  
If not, can one construct or characterize non-Sasakian critical points, and how would their Tanaka–Webster torsion or transverse geometry behave?

Why is it a "good" problem:  
This problem is carefully positioned at the meeting point of **contact geometry**, **CR geometry**, and **geometric analysis**. It correctly introduces a bona fide infinite-dimensional deformation space—the set of contact forms compatible with a fixed CR structure—and poses a variational question whose Euler–Lagrange characterization is unknown and nontrivial from both analytic and geometric viewpoints.  

It forces one to compare **Riemannian** variational principles (total scalar curvature extremality) with the known **transverse** or **Tanaka–Webster** frameworks, potentially unifying them or revealing a distinct mechanism selecting Sasakian structures. Conceptually simple (“are Sasakian metrics variationally characterized within their CR deformation class?”) yet analytically and geometrically profound, this problem could open a new variational bridge between CR geometry and global Riemannian curvature functionals, analogous to the Calabi extremal theory in Kähler geometry but in the non-integrable contact world.